\documentclass[11pt]{article}
\usepackage{icmlko}
\begin{document}

\title{모바일만 반대로 기댄다 --- 여덟 노트째 남은 마지막 이상을 닫는다}
\author{월드모델 실험실}
\date{노트 487 · 2026년 8월 3일}
\maketitle

\begin{abstract}
노트 $481 \cdot 486$ 이 모바일 하나를 이상으로 남겼다 --- \emph{대조
격차($+0.2497$)가 날 격차($+0.2244$)보다 크다.} 덮음이 샌 것도(노트 481
이 고침) 라벨 압축도(노트 486) 아니었고, 마지막 용의자는
\textbf{대조 팔의 재학습}이었다.

\textbf{반증됐다.} 재학습을 없앤 대조(챔피언 그대로 두고 예보 때만 가림)
에서 모바일은 $+0.2559$ --- \textbf{여전히 날 격차보다 크다}.

그런데 그 표가 다른 것을 보여 줬다. \textbf{네 도메인은 두 설계가 거의
같고}(차 $\le 0.017$) \textbf{게임만 $0.137$ 로 흔들린다}($+0.0018$ 대
$-0.1356$).

그래서 격차 말고 \textbf{두 절반의 값}을 봤다.

\begin{center}
\begin{tabular}{lrr}
\hline
저덮음 축을 가리면 & 축 많은 절반 & 축 적은 절반 \\
\hline
게임 & $-0.2623$ & $-0.0086$ \\
만화 & $-0.1511$ & $-0.0003$ \\
세계애니 & $-0.0886$ & $-0.0353$ \\
애니 & $-0.0868$ & $-0.0274$ \\
\textbf{모바일} & $\mathbf{+0.0167}$ & $\mathbf{-0.0148}$ \\
\hline
\end{tabular}
\end{center}

\textbf{넷은 당연한 방향이다} --- 많은 쪽이 잃을 축이 더 많으니 더 잃는다.
\textbf{모바일만 부호가 갈린다}: 가리면 많은 쪽이 \emph{오르고} 적은 쪽이
내린다.
\end{abstract}

\begin{figure}[t]
\centering
\includegraphics[width=\textwidth]{figs/f1.pdf}
\caption{\textbf{가} --- 다섯 도메인. 모바일만 파랑이 $0$ 위에 있다.
\textbf{나} --- 모바일 두 절반.}
\end{figure}

\section{이상의 정체}

\textbf{모바일의 저덮음 축들은 축 많은 행에 해롭고 적은 행에 이롭다.}
가리면 많은 쪽이 $+0.0167$ 오르고 적은 쪽이 $-0.0148$ 내린다 --- 전체로는
$\approx 0$ 이지만 \emph{격차는 벌어진다}. 그것이 노트 481 부터 여덟
노트째 ``이상''으로 보이던 것의 전부다.

노트 441 의 무늬가 여기서 다시 나온다 --- \emph{어떤 축은 어떤 행에서
거꾸로 쓰인다.} 다만 그때는 \textbf{수집 단위}였고 여기는 \textbf{도메인
안의 부분 무리} 단위다.

\textbf{채택거리는 아니다.} 모바일 전체로는 $\approx 0$ 이고, 판 무게
$441/3{,}430$ 을 곱하면 더 작다. \textbf{설명이지 개선이 아니다.}

\section{게임은 설계에 민감하다}

부수적으로 하나 더 나왔다. \textbf{게임만 두 대조 설계가 크게 다르다}
--- 재학습 대조 $+0.0018$, 재학습 없는 대조 $-0.1356$.

게임은 자기만 갖는 축이 넷이다(가격 · 연령등급 · 램 · 카테고리 수).
\emph{다시 학습하면} 모형이 그 축들 없이도 쓸 만한 대체를 배우지만,
\emph{예보 때만 빼면} 그 축들에 기대던 예보가 통째로 무너진다. 축 많은
절반이 $-0.2623$ 로 제일 크게 잃는 것이 그 증거다.

\begin{quote}
\textbf{조항} --- ``이 자료가 얼마나 값어치가 있나''를 잴 때 \textbf{재학습
대조와 재학습 없는 대조는 다른 물음에 답한다.} 앞엣것은 \emph{그 자료가
없는 세상}을, 뒤엣것은 \emph{그 자료를 지금 잃는 것}을 잰다. 어느 쪽을
쟀는지 적는다.
\end{quote}

\section*{여덟 노트를 닫으며}

노트 $480 \sim 487$ 이 하나의 물음을 쫓았다 --- \emph{축이 많은 행과 적은
행의 격차는 무엇인가.} 답은 여러 조각이었다.

\begin{itemize}
\item \textbf{행 난이도}(노트 480) --- 맞지만 재는 변수를 못 댔다(노트 483).
\item \textbf{라벨 압축}(노트 $484 \cdot 486$) --- 애니의 $89\%$ 를
설명하고 나머지 넷에서는 반대로 가린다.
\item \textbf{축 자체}(노트 481) --- 도메인마다 $+98\% \sim -67\%$ 로
퍼져 평균이 못 쓸 요약이다.
\item \textbf{축이 거꾸로 쓰이는 무리}(이 노트) --- 모바일 하나.
\end{itemize}

\textbf{한 물음에 답이 넷이고 도메인마다 무게가 다르다.} 이것이 이
자료에서 ``왜''를 묻는 일의 값이자 한계다.

\section*{남은 것}

\begin{itemize}
\item \textbf{모바일의 어느 축이 거꾸로인가} --- 저덮음 축 전체를 한꺼번에
가렸다. 하나씩 갈라야 노트 441 처럼 이름을 댈 수 있다.
\item \textbf{만화는 압축이 같은 방향인 유일한 도메인}(노트 486) --- 아직
안 봤다.
\item \textbf{게임의 전용 축 넷} --- 재학습 여부로 $0.137$ 이 갈리는
크기다. 그 축들이 실제로 무엇을 하는지 직접 안 봤다.
\end{itemize}

\end{document}
